\section{Experiments}
We separate two questions: does training on expanded graphs improve the simulator, and does residual risk identify a better placement of a fixed expansion budget? Fresh Goop training and paired WaterDrop continuations test both. Earlier WaterDrop controls examine why error ranking can disagree with action value. The original Sand and WaterDrop comparisons remain in Appendix~\ref{sec:historical-audit}.

\subsection{Learning to use expanded graphs}
We train six Goop models from initialization for 100,000 updates: three paired seeds each for base-only and mixed-graph training. All use faithful regression, width 128, ten message-passing blocks and all 1,000 official training trajectories. The mixed arm adds a uniformly random quarter of the optional pairs independently with probability $1/2$. Initial weights, sampled frames and noise are paired by seed; the native capped base graph and self-messages are preserved. At evaluation, random25, speed25, risk25 and relative-velocity-RMS25 use the same quarter-annulus budget, with radius ratio 1.267. Base and dense control the amount of expansion.

Graph exposure improves full-rollout error under an unchanged random25 evaluation policy: mean H395 position MSE falls from $0.301\pm0.169$ to $0.129\pm0.031$, and every paired seed improves. The observed-state comparison separates that training effect from placement (Table~\ref{tab:goop-exposure-placement}). Previous-observed risk remains worse than random in every mixed-model seed, despite a smaller deficit than with base-only training. Its mean frame Spearman correlation is $.214\pm.088$ with base error but $.001\pm.048$ with its own sparse-action benefit. Exposure therefore improves this simulator without making residual rank a reliable guide to useful extra messages.

\begin{table}[t]
\centering\small
\caption{Observed-test position-MSE contrasts, $\times10^{-10}$, within each dataset. Goop trains from initialization to 100k; WaterDrop continues three 100k parents to 110k. Values are means and sample SDs of three paired seeds; negative values favor the first term. The interaction changes the risk-minus-random gap after exposure. These scales do not compare effect size across materials.}
\label{tab:goop-exposure-placement}
\setlength{\tabcolsep}{3pt}
\begin{tabular}{lrr}
\toprule
Contrast & Goop & WaterDrop\\
\midrule
Mix $-$ base, base policy & $-3.04\pm5.12$ & $+0.70\pm0.40$\\
Mix $-$ base, random25 & $-5.51\pm4.28$ & $+0.28\pm0.92$\\
Risk $-$ random, base train. & $+6.39\pm1.45$ & $+1.62\pm1.07$\\
Risk $-$ random, mixed train. & $+2.84\pm3.62$ & $-0.99\pm0.67$\\
Training interaction & $-3.55\pm2.18$ & $-2.61\pm0.74$\\
\bottomrule
\end{tabular}
\end{table}

Observed-state evaluation uses 150 histories per model and recomputes risk on the preceding observed base history. Autonomous cached risk instead reuses the score from its own preceding selected graph. We average frames within trajectories, then trajectories within each seed; all three seed values enter the reported means and sample SDs. Figure~\ref{fig:goop-exposure-placement} in the appendix shows every paired seed. The Goop study follows the WaterDrop analyses and is exploratory; its endpoints and policies were fixed before evaluation.

A separate WaterDrop follow-up pairs 10k base-only and mixed-graph continuations from three faithful 100k parents, inheriting optimizer state and matched frame/noise schedules. Its exposure effect is conditional on those parents, distinct from Goop training from initialization. Mixed training changes the observed-test risk-minus-random position-MSE gap from $(+1.62\pm1.07)\times10^{-10}$ to $(-0.99\pm0.67)\times10^{-10}$, with every seed changing sign (Table~\ref{tab:goop-exposure-placement}). Validation's relative gap also improves, but mixed risk remains worse than random on average. All 810 H995 autonomous outcomes complete. Under random25, mix-minus-base MSE is $-.0060\pm.0056$, improving in every seed; dense, speed and cached risk also improve in each seed. Yet the autonomous risk-minus-random interaction is $+.00284\pm.00466$, with mixed seed signs. Exposure can therefore improve placement on observed test histories without a consistent autonomous placement gain. Substantial box excursions and increased recorded rollout times remain (Appendix~\ref{sec:waterdrop-110k-exposure}).

\subsection{Does the effect survive autonomous feedback?}
All six policies are evaluated for all 395 forecasts on 30 Goop test trajectories per model. Of 1,080 required outcomes, 1,077 complete and three hit the candidate-pair guard, all in mixed seed 2: base, dense and cached risk each fail once. Their three-seed full-horizon means remain undefined (Table~\ref{tab:goop-rollouts}), as does the full-horizon risk-minus-random training interaction. The cached-risk comparison with random changes sign across the two defined seeds; we do not average those survivors. Pointwise MSE at forecast 200 remains a secondary endpoint and does not replace the full horizon.

\begin{table}[t]
\centering\small
\caption{Goop autonomous position MSE. H395 averages all 395 forecast errors; @200 is one forecast. Each row requires 90 outcomes. A dash denotes an undefined three-seed statistic after a failed required outcome. No outcome is missing.}
\label{tab:goop-rollouts}
\setlength{\tabcolsep}{3pt}
\begin{tabular}{lrrr}
\toprule
Training / policy & H395 & @200 & Failed\\
\midrule
Base / base & $.247\pm.149$ & $.249\pm.094$ & 0\\
Base / random25 & $.301\pm.169$ & $.300\pm.098$ & 0\\
Mix / random25 & $.129\pm.031$ & $.148\pm.013$ & 0\\
Mix / cached risk25 & --- & $.155\pm.023$ & 1\\
\bottomrule
\end{tabular}
\end{table}

Among complete mixed-training policies, speed25 and relative-velocity-RMS25 have full-horizon MSE $.112\pm.027$ and $.125\pm.032$ and beat random25 in every seed. Accuracy gains still leave substantial geometric error: mixed random25, speed25 and RMS25 place $22.16\%$, $19.55\%$ and $20.53\%$ of particles beyond the metadata box by more than $10^{-6}$, versus $0.0632\%$ for matching truth. Their mean trajectory-maximum excursions are $.401$, $.417$ and $.426$, versus $.000722$ for truth. These measure box excursions, not conservation or certified physical failures. The appendix retains every policy and seed, all failures, and these boundary and cost diagnostics (Appendix~\ref{sec:goop-graph-exposure}). A fixed metadata-selected rollout illustrates the remaining physical mismatch in Figure~\ref{fig:qualitative-goop2d}.

\subsection{Why can residual ranking misallocate?}
The earlier WaterDrop controls isolate residual ranking from action value: at 100k updates, faithful/NLL risk correlates with base error ($.357\pm.060$ / $.411\pm.023$), but weakly with actual sparse benefit ($-.021\pm.049$ / $-.052\pm.013$); random has lower one-step error in every seed. Restoring trained self-messages reduces error without reversing that validation/test ordering. Risk's larger alignment gain is outweighed by its squared prediction change in every full-model seed. Autonomous signs vary, with failures and box excursions under both graph conventions (Appendices~\ref{sec:compact-study}, \ref{sec:actual-action-benefit}, \ref{sec:graph-bridge}, \ref{sec:optional-exposure}, \ref{sec:full-results} and~\ref{sec:native-rollout-followup}). These findings motivate testing exposure separately from placement.

\subsection{What computation does the policy require?}
Previous-observed risk includes an extra scoring pass; autonomous cached risk reuses its previous forward output. Their timings measure different computations. Fixed-order autonomous measurements also follow different geometries, so elapsed time and edge counts do not establish a causal speedup. The appendix retains measured cost alongside every policy's error and failure counts. The studies use only three training seeds, reuse earlier evidence to develop follow-up questions, and remain distinct from the original saved-model comparisons.
